% Chapter 3 -- Methodology (\input by main.tex): detection pipeline + object completion
% (completion half merged in below, formerly sec_6_completion_method.tex).
%
% Honesty-phrasing rules applied (personal_agenda P1):
%   - production classifier is the FROM-SCRATCH Stage B checkpoint
%     (stage_b_scratch_best.pth, Finding #31); synthetic Stage A pretraining
%     was DROPPED from the final pipeline and is kept only as Ch-4 ablation
%     material (#25/#30). We do NOT present two-stage training as a contribution.
%   - measured vs inferred evidence tiers marked in prose; timing numbers are
%     measured (timing_seq08_full_n4068.json), cited as such.
%   - every number traceable to a finding / output artifact.
%   - "the classifier's job is precision" claim deferred to the Ch-4 ablation.
%
\chapter{Proposed Methodology}
\label{ch:pipeline}

This chapter presents the proposed method: a modular detection-and-completion system
for automotive LiDAR, and the donor-frame metric that makes its completion measurable
on real data. Section~\ref{sec:meth-architecture} gives the overall architecture; the
sections that follow describe each stage in the order a scan flows through it, and the
completion is evaluated with the donor metric in Section~\ref{sec:donor-validation}.

\section{System Architecture}
\label{sec:meth-architecture}

The method is a single system that turns a raw LiDAR scan into completed,
box-estimated cars together with a quality measurement for each completion. It is
organised as three stages in sequence, each the subject of one section below. The
\emph{detection} stage (Section~\ref{sec:meth-detection}) segments and classifies car
clusters from the raw scan, emitting an oriented bounding box and a stable track
identity per car. The \emph{completion} stage (Section~\ref{ch:completion-method})
reconstructs each detected car's full surface from the one-sided partial a single
frame observes. The \emph{measurement} stage --- the donor-frame coverage metric
(Section~\ref{sec:donor-metric}) --- makes that completion quantifiable on real data,
where no clean ground-truth complete shape exists. Figure~\ref{fig:pipeline} shows the
detection chain that produces every input the later two stages consume, annotated with
the measured per-stage cost.

The design is deliberately modular and classical: classical geometry does the work
geometry solves well, and small learned models are used only where discrimination
needs them, so every stage is separately inspectable and its cost separately
measurable. This architecture is what lets the evaluation
(Chapter~\ref{ch:evaluation}) attribute the recall limit to a single stage rather than
to an opaque network.

\section{Detection Pipeline}
\label{sec:meth-detection}

This section describes the car-detection pipeline that produces every detection
evaluated in Chapter~\ref{ch:evaluation} and every cluster completed in
Section~\ref{ch:completion-method} and Section~\ref{sec:donor-validation}. The
design is deliberately modular and classical: a chain of independent stages,
each with a small number of interpretable parameters, rather than an end-to-end
learned detector. This choice trades peak accuracy for interpretability,
reproducibility, and near-zero manual annotation --- the classifier is trained on
labels mined automatically from the pipeline itself
(Section~\ref{sec:pipeline-classifier}) --- and it is what lets Section~\ref{sec:recall-analysis}
root-cause the recall shortfall to a single stage rather than to an opaque
network.

The pipeline takes one raw LiDAR scan and emits a set of car detections, each an
oriented bounding box with a stable track identity across frames. Every stage
parameter lives in a single dictionary, \texttt{PIPELINE\_CONFIG} in
\texttt{src/pipeline.py}, which was frozen before the write-up began; the frozen
values are collected in Table~\ref{tab:pipeline-config}. Figure~\ref{fig:pipeline}
gives the full chain with the measured per-stage cost.

\subsection{Overview of the Modular Detection Pipeline}
\label{sec:pipeline-overview}

\begin{figure}[H]
  \centering
  \begin{tikzpicture}[
      font=\small,
      node distance=6mm,
      stage/.style={
        draw, rounded corners, align=center,
        minimum width=34mm, minimum height=11mm,
        fill=blue!5, thick},
      learned/.style={stage, fill=orange!12},
      io/.style={draw, ellipse, align=center, fill=black!5,
        minimum width=28mm, minimum height=9mm},
      ms/.style={font=\footnotesize\itshape, text=black!60},
      arr/.style={-{Stealth[length=2.2mm]}, thick},
      grp/.style={draw=black!45, dashed, rounded corners, inner sep=3mm},
      grplab/.style={font=\footnotesize\bfseries, text=black!65},
    ]
    \node[io] (raw) {Raw LiDAR scan\\$\sim$120k points};
    \node[stage, below=of raw] (pre) {Preprocess\\{\footnotesize $z$-filter $\to$ voxel $\to$ denoise}};
    \node[stage, below=of pre] (ground) {Ground removal\\{\footnotesize RANSAC plane}};
    \node[stage, below=of ground] (cluster) {Clustering\\{\footnotesize HDBSCAN @ \SI{0.10}{\meter}}};
    \node[stage, below=12mm of cluster] (geo) {Geometric filter\\{\footnotesize OBB gates}};
    \node[learned, below=of geo] (clf) {Classifier\\{\footnotesize PointNet, car/not-car}};
    \node[stage, below=of clf] (track) {Tracker\\{\footnotesize centroid + track vote}};
    \node[io, below=of track] (out) {Car detections\\{\footnotesize boxes + track ids}};
    \node[learned, right=24mm of out] (comp) {Completion\\{\footnotesize PCN (\S\ref{ch:completion-method})}};

    \draw[arr] (raw) -- (pre);
    \draw[arr] (pre) -- (ground);
    \draw[arr] (ground) -- (cluster);
    \draw[arr] (cluster) -- (geo);
    \draw[arr] (geo) -- (clf);
    \draw[arr] (clf) -- (track);
    \draw[arr] (track) -- (out);
    \draw[arr] (out) -- (comp);

    % dashed phase blocks (segmentation: points -> clusters; detection: clusters -> cars)
    \begin{scope}[on background layer]
      \node[grp, fit=(pre)(ground)(cluster)] (segbox) {};
      \node[grp, fit=(geo)(clf)(track)] (detbox) {};
    \end{scope}
    \node[grplab, rotate=90] at ([xshift=-4mm]segbox.west) {Segmentation};
    \node[grplab, rotate=90] at ([xshift=-4mm]detbox.west) {Detection};

    % per-stage timing in one column, anchored to the phase-block right edge (uniform gap)
    \node[ms, anchor=west] at ([xshift=6mm]segbox.east|-pre)     {\SI{126}{\milli\second} (20\%)};
    \node[ms, anchor=west] at ([xshift=6mm]segbox.east|-ground)  {\SI{56}{\milli\second} (9\%)};
    \node[ms, anchor=west] at ([xshift=6mm]segbox.east|-cluster) {\SI{370}{\milli\second} (59\%)};
    \node[ms, anchor=west] at ([xshift=6mm]detbox.east|-geo)     {\SI{36}{\milli\second} (6\%)};
    \node[ms, anchor=west] at ([xshift=6mm]detbox.east|-clf)     {\SI{34}{\milli\second} (5\%)};
    \node[ms, anchor=west] at ([xshift=6mm]detbox.east|-track)   {\SI{0.3}{\milli\second}};
    \node[ms, below=1mm of comp.south, text width=34mm, align=center]
      {\footnotesize\itshape +\SI{13}{\milli\second} per completed car};
  \end{tikzpicture}
  \caption[Detection pipeline]{The detection pipeline, grouped into two phases (dashed): a
    \emph{segmentation} phase that turns a raw scan into clusters (preprocess, ground removal,
    clustering) and a \emph{detection} phase that turns clusters into tracked car boxes
    (geometric filter, classifier, tracker); the optional completion stage (right) runs off the
    tracked detections. Learned stages are shaded; all others are
    classical. Per-stage figures are mean wall-clock time on the full
    sequence~08 under the final (promoted) configuration
    (\num{4068} frames, AMD Ryzen 7 7800X3D + RTX 3070~Ti;
    \texttt{output/experiments/timing/timing\_seq08\_full\_n4068\_combined.json}), and are
    measured, not estimated. The total is
    \SI{623.5}{\milli\second\per\frame} ($\approx$\SI{1.6}{\fps}); the only learned
    \emph{per-frame} stage --- the classifier --- accounts for just
    \SI{34}{\milli\second} (\SI{5}{\percent}), and the classical clustering and
    ground-removal stages dominate. Completion runs off the tracked detections and
    is optional; it is learned, but its cost is per completed car, not per frame, so
    it does not enter this per-frame budget.}
  \label{fig:pipeline}
\end{figure}

The system is \emph{offline} by design: the tracker and the track-level filter
use the whole sequence, and real-time operation is out of scope. At
\SI{623.5}{\milli\second\per\frame} the pipeline is roughly \SI{1.6}{\fps}, an order
of magnitude below sensor rate; we characterise this cost honestly in the runtime
section of Chapter~\ref{ch:evaluation} rather than claiming real-time capability. The cost lives in the classical stages:
clustering alone is \SI{59}{\percent} of the frame budget, and the only learned
per-frame stage, the classifier, is just \SI{5}{\percent}. The expensive parts of this pipeline are
classical, not neural.

\subsection{Preprocessing: Filtering, Downsampling, and Denoising}
\label{sec:pipeline-preprocess}

Each raw scan (roughly \num{120000} points) passes through three cleaning steps
before any object reasoning:

\begin{enumerate}
  \item \textbf{Ground-height $z$-filter.} Points below
    $z = \SI{-2.0}{\meter}$ in the sensor frame are dropped. This removes the bulk
    of the road surface and any sub-ground noise cheaply, before the more
    expensive steps see them.
  \item \textbf{Voxel downsampling} at \SI{0.05}{\meter}. Points are collapsed to
    one representative per \SI{5}{\centi\meter} voxel, which both regularises the
    range-dependent density of LiDAR and cuts the point count by roughly an order
    of magnitude.
  \item \textbf{Statistical outlier removal.} Points whose mean distance to their
    \num{20} nearest neighbours exceeds \num{2.0} standard deviations of the
    local distribution are discarded, removing isolated speckle.
\end{enumerate}

The order of the last two steps is behaviour-changing. In the frozen configuration
(\texttt{voxel\_\allowbreak before\_\allowbreak denoise = True}) downsampling runs \emph{before}
denoising, so the statistical outlier test operates on roughly ten times fewer
points --- and on a regularised point set, which changes which points survive.
This ordering was adopted as part of the pipeline optimisation primarily
because it improves detection (discussed in Chapter~\ref{ch:evaluation}); its effect on the measured
denoising cost is small. It is behaviour-changing, so we flag it here rather than
presenting the two orderings as equivalent.
Preprocessing costs \SI{126}{\milli\second\per\frame} in total, split across the
$z$-filter (\SI{58}{\milli\second}), denoising (\SI{50}{\milli\second}), and
downsampling (\SI{18}{\milli\second}).

\subsection{Ground-Plane Removal by RANSAC}
\label{sec:pipeline-ground}

The dominant remaining structure after preprocessing is the ground plane, which
must be removed before clustering or it bridges every object into one component.
We fit a single plane with RANSAC: over \num{300} iterations, three points are
sampled, their plane computed, and inliers counted within
\SI{0.2}{\meter}; the best-supported plane is kept. The plane is
\emph{accepted only if its normal is near-vertical} --- the $z$-component of the
unit normal must exceed \num{0.5} --- so a spuriously-fitted vertical surface (a
wall, the side of a truck) cannot be mistaken for ground. If the check fails, no
points are removed and the whole cloud is passed on as objects; this fail-safe is
rare but prevents a bad plane fit from deleting real cars.

A single global plane is a deliberate simplification: it assumes the local ground
is approximately planar, which holds on the flat urban sequences used here but
would degrade on strong slopes. A grid-based ground model (Patchwork++-style) is
listed as future work, not a limitation we claim to have solved. The iteration
count was reduced from \num{1000} to \num{300} as part of the runtime optimisation after
confirming the plane fit converges well before \num{1000} iterations on this
data; the fit uses a fixed seed (\texttt{np.random.default\_rng(42)}) so ground
removal, and therefore the whole pipeline, is deterministic across runs. Ground
removal costs \SI{56}{\milli\second\per\frame} (\SI{9}{\percent} of the frame).

\subsection{Object-Candidate Clustering with HDBSCAN}
\label{sec:pipeline-cluster}

Non-ground points are grouped into object candidates with HDBSCAN, a
density-based hierarchical clusterer, using a minimum cluster size of \num{10}
and minimum samples of \num{5}. HDBSCAN is chosen over a fixed-radius method
because LiDAR density falls off with range: a single Euclidean radius that merges
nothing at \SI{5}{\meter} will over-segment at \SI{40}{\meter}. HDBSCAN's
variable-density hierarchy adapts to this without a hand-tuned radius schedule,
and Chapter~\ref{ch:evaluation} reports a controlled benchmark confirming it beats
DBSCAN and Euclidean extraction on recall at matched settings.

Two details of the frozen configuration are load-bearing. First, clustering runs
on a \emph{coarser} \SI{0.10}{\meter} voxel grid than the \SI{0.05}{\meter}
preprocessing grid: the object cloud is re-downsampled to \SI{0.10}{\meter},
HDBSCAN is run there, and the labels are propagated back to the full-resolution
points by nearest neighbour. This was adopted because the coarser
grid closes intra-car density gaps that otherwise cause a single car to split
into several clusters --- the central mechanism of the recall bottleneck analysed
in Section~\ref{sec:recall-analysis} (Figure~\ref{fig:clustering}) --- and it added \num{1655} true positives on the full sequence~08.
Second, HDBSCAN is by far the most expensive stage at
\SI{370}{\milli\second\per\frame} (\SI{59}{\percent} of the budget), because it is
superlinear in the point count; the coarse-grid trick is as much a speed measure
as an accuracy one. The clustering stage produces on average \num{45} candidate
clusters per frame.

\begin{figure}[H]
  \centering
  \begin{tikzpicture}[
      font=\footnotesize, node distance=5mm and 7mm,
      blk/.style={draw, rounded corners, align=center, fill=blue!5, thick,
        minimum height=9mm, inner sep=3pt},
      io/.style={draw, ellipse, align=center, fill=black!5, inner sep=2pt},
      arr/.style={-{Stealth[length=2mm]}, thick},
      note/.style={font=\scriptsize\itshape, text=black!60, align=center},
    ]
    % --- top: the clustering flow ---
    \node[io] (in) {non-ground\\points};
    \node[blk, right=of in] (vox) {re-voxel\\\SI{0.10}{\meter}};
    \node[blk, right=of vox] (hdb) {HDBSCAN\\{\scriptsize min\_size \num{10}, min\_samples \num{5}}};
    \node[blk, right=of hdb] (prop) {propagate labels\\{\scriptsize NN $\to$ full-res}};
    \node[io, right=of prop] (out) {$\sim$\num{45}\\clusters};
    \draw[arr] (in)--(vox); \draw[arr] (vox)--(hdb);
    \draw[arr] (hdb)--(prop); \draw[arr] (prop)--(out);

    % --- bottom: why the coarser grid (split vs whole) ---
    \node[note] (why) at ($(hdb.south)+(0,-6mm)$)
      {\textbf{Why re-voxel coarser?} closing intra-car density gaps ($+\num{1655}$ TP on seq~08)};
    \begin{scope}[shift={($(vox.south)+(0,-17mm)$)}]
      \draw[rounded corners, black!45] (-1.15,-0.42) rectangle (1.15,0.42);
      \fill[blue!30] (-0.62,0) ellipse (0.42 and 0.30);
      \fill[red!30]  (0.62,0)  ellipse (0.42 and 0.30);
      \node[note] at (0,-0.85) {\SI{0.05}{\meter} grid:\\car splits (2 clusters) {\textcolor{red!75!black}{$\times$}}};
    \end{scope}
    \begin{scope}[shift={($(prop.south)+(0,-17mm)$)}]
      \draw[rounded corners, black!45] (-1.15,-0.42) rectangle (1.15,0.42);
      \fill[green!40] (0,0) ellipse (1.0 and 0.30);
      \node[note] at (0,-0.85) {\SI{0.10}{\meter} grid:\\one cluster {\textcolor{green!55!black}{$\surd$}}};
    \end{scope}
  \end{tikzpicture}
  \caption[Clustering stage]{The clustering stage. Non-ground points are re-voxelised to a coarser
    \SI{0.10}{\meter} grid, clustered with HDBSCAN (density-adaptive, so a single car stays whole at
    both near and far range without a hand-tuned radius), and the cluster labels are propagated back to
    the full-resolution points by nearest neighbour. The coarser grid is load-bearing (bottom): on the
    \SI{0.05}{\meter} preprocessing grid intra-car density gaps split one car into several clusters,
    whereas the \SI{0.10}{\meter} grid closes them and keeps the car a single cluster --- worth
    \num{1655} true positives on sequence~08. HDBSCAN is the pipeline's most expensive stage at
    \SI{370}{\milli\second\per\frame} (\SI{59}{\percent}), so the coarse grid is also a speed measure.}
  \label{fig:clustering}
\end{figure}

\subsection{Geometric Filtering of Non-Car Clusters}
\label{sec:pipeline-geometric}

Most clusters are not cars: they are trees, poles, building facades, and
fragments. A geometric filter rejects clusters whose oriented bounding box (OBB)
is inconsistent with a car, using only cheap, interpretable shape gates. For each
cluster the OBB is computed and the cluster is kept only if it passes \emph{all}
of the following:

\begin{itemize}
  \item at least \num{15} points;
  \item OBB volume in $[\num{0.5}, \num{50}]\,\si{\meter\cubed}$;
  \item longest OBB extent $\leq \SI{6.0}{\meter}$, and the two largest extents
    above small floors (\SI{0.5}{} and \SI{0.4}{\meter}), rejecting slivers;
  \item box centre no more than \SI{1.5}{\meter} above the fitted ground plane,
    and vertical point span $\leq \SI{1.8}{\meter}$, rejecting tall structures;
  \item aspect ratio (longest\,/\,shortest extent) $\leq \num{6.0}$, rejecting
    thin walls and poles.
\end{itemize}

These bounds are generous car priors, not a tight car model --- they are meant to
remove obvious non-cars while leaving the fine car/not-car decision to the
learned classifier. The height gate is computed \emph{relative to the fitted
ground plane}, not to a fixed world height, so it is robust to the sensor not
being perfectly level. The filter costs \SI{36}{\milli\second\per\frame}
(\SI{6}{\percent}). Its individual gates are ablated in Chapter~\ref{ch:evaluation}; the stage is
high-recall and low-precision by construction, and the
classifier is what turns its survivors into precise detections.

\subsection{The Car / Not-Car Point Classifier}
\label{sec:pipeline-classifier}

Each surviving cluster is classified as \emph{car} or \emph{not-car} by a small
dual-branch PointNet (\texttt{src/\allowbreak classifier.py}). The binary decision
replaces
an earlier multi-class (car/bus/motorcycle) design that was dropped once the data
showed near-zero real bus and motorcycle instances, leaving the car-only detector
(Chapter~\ref{ch:introduction}).

\paragraph{Architecture.} The network has two branches
(Figure~\ref{fig:classifier}) and about \num{80000}
parameters --- deliberately small:
\begin{itemize}
  \item a \textbf{point branch} that takes \num{512} points (sampled or padded)
    normalised into the unit sphere, and applies shared per-point MLPs
    ($3 \to 64 \to 128 \to 256$) followed by a max-pool to a \num{256}-d shape
    descriptor;
  \item a \textbf{bounding-box branch} that takes an \num{8}-dimensional metric
    feature vector (sorted OBB extents, volume, two aspect ratios,
    $\log(1+\text{point count})$, vertical span), $z$-score normalised, through a
    small MLP to \num{32}-d.
\end{itemize}
The two are concatenated and passed through a small head
($288 \to 128 \to 2$) with dropout. The bounding-box branch is what re-injects
absolute metric scale, which the unit-sphere normalisation deliberately strips
from the point branch: a distant small car and a nearby one have the same
normalised shape but different real extents, and the metric branch is how the
network sees the difference. There are no T-Nets, since clusters already sit in a
consistent sensor frame.

\begin{figure}[H]
  \centering
  \resizebox{\linewidth}{!}{%
  \begin{tikzpicture}[
      font=\footnotesize, node distance=6mm and 7mm,
      blk/.style={draw, rounded corners, align=center, fill=blue!5, thick,
        minimum height=8mm, inner sep=3pt},
      lrn/.style={blk, fill=orange!12},
      io/.style={draw, ellipse, align=center, fill=black!5, inner sep=2pt},
      arr/.style={-{Stealth[length=2mm]}, thick},
    ]
    \node[io] (cl) at (0,0) {cluster};
    \node[blk] (pts) at (2.3,1.3) {\num{512} points\\unit-sphere};
    \node[lrn, right=of pts] (pmlp) {shared MLP\\{\scriptsize $3\!\to\!64\!\to\!128\!\to\!256$}};
    \node[blk, right=of pmlp] (pool) {max-\\pool};
    \node[blk, right=of pool] (pd) {shape\\\num{256}-d};
    \node[blk] (feat) at (2.3,-1.3) {\num{8} OBB\\features};
    \node[lrn, right=of feat] (bmlp) {MLP};
    \node[blk, right=of bmlp] (bd) {metric\\\num{32}-d};
    \node[blk] (cat) at ($(pd.east)+(1.1,-1.3)$) {concat\\\num{288}-d};
    \node[lrn, right=of cat] (head) {head\\{\scriptsize $288\!\to\!128\!\to\!2$}};
    \node[io, right=of head] (out) {car /\\not-car};
    \draw[arr] (cl) |- (pts.west);
    \draw[arr] (cl) |- (feat.west);
    \draw[arr] (pts)--(pmlp); \draw[arr] (pmlp)--(pool); \draw[arr] (pool)--(pd);
    \draw[arr] (feat)--(bmlp); \draw[arr] (bmlp)--(bd);
    \draw[arr] (pd) -| (cat.north);
    \draw[arr] (bd) -| (cat.south);
    \draw[arr] (cat) -- (head);
    \draw[arr] (head) -- (out);
  \end{tikzpicture}%
  }
  \caption[Dual-branch classifier]{The dual-branch classifier (\texttt{src/classifier.py}), applied to one
    cluster. The \emph{point branch} takes \num{512} unit-sphere-normalised points
    through shared per-point MLPs and a max-pool to a \num{256}-d shape descriptor;
    the \emph{bounding-box branch} takes \num{8} metric shape features --- the
    absolute scale the sphere normalisation strips --- to \num{32}-d. The two are
    concatenated and passed through a small head to a car/not-car decision. Learned
    blocks are shaded.}
  \label{fig:classifier}
\end{figure}

\paragraph{Training data (mined, not hand-labelled).} The classifier is trained
on real LiDAR clusters mined automatically from SemanticKITTI
(\texttt{src/mine\_stage\_b.py}). For each training frame the pipeline's own
stages~1--5 are replayed, the ground-truth semantic labels are propagated onto the
surviving clusters by nearest neighbour, and each cluster is labelled \emph{car}
or \emph{not-car} by majority vote, keeping only clusters whose majority class
has purity $\geq \num{0.75}$. This yields \num{420333} training clusters
(sequences 00--07, 09, 10) and \num{130394} validation clusters (sequence~08),
with the intrinsic $\approx$21\,\%/79\,\% car/not-car imbalance handled by a
class-weighted loss. The key property is that the training clusters are produced
by the \emph{same} clustering and filtering the classifier meets at inference ---
including HDBSCAN splits and filter survivors --- so it learns on exactly the
distribution it must classify, at no manual-annotation cost.

\paragraph{Production checkpoint: trained from scratch.} The production classifier
is \texttt{stage\_b\_\allowbreak scratch\_\allowbreak best.pth}, trained on the
mined real clusters from
random initialisation (val macro-$F_1$ \num{0.9285}). An earlier
two-stage design first pretrained on synthetic ShapeNet partial renders
(``Stage~A'') and then fine-tuned on the real clusters (``Stage~B''). We report
the two-stage machinery in Chapter~\ref{ch:evaluation} as an \emph{ablation}, because a controlled
comparison showed the synthetic pretraining to be \emph{redundant}
given \num{420000} real clusters: the from-scratch classifier matches the
pretrained one (macro-$F_1$ \num{0.9285} vs.\ \num{0.9225}). The final pipeline
therefore uses the simpler from-scratch checkpoint, and the synthetic
pretraining, together with the finding that the sim-to-real gap is total and
symmetric, is presented as a negative result rather than as a
contribution. The classifier costs \SI{34}{\milli\second\per\frame} (about
\SI{0.75}{\milli\second} per cluster over the \num{45} clusters in a frame).

\paragraph{Training.} The from-scratch checkpoint is trained for \num{15} epochs
with Adam (learning rate \num{e-4}, weight decay \num{e-5}) and a step schedule
that halves the rate every \num{10} epochs, at batch size \num{64}. The loss is
cross-entropy weighted by inverse square-root class frequency (normalised to unit
mean), which offsets the intrinsic $\approx$21\,\%/79\,\% car/not-car imbalance
without resampling. Training clusters are augmented online with a random rotation
about the vertical axis and Gaussian jitter ($\sigma = \SI{0.005}{\meter}$); the
input is the \num{512}-point cluster of the architecture above. Runs are seeded
and the reported checkpoint is the epoch with the best validation macro-$F_1$
(\num{0.9285} on the sequence-08 validation clusters).

\subsection{Multi-Frame Tracking and Track-Level Filtering}
\label{sec:pipeline-tracker}

Detections are linked across frames by a centroid tracker
(\texttt{src/tracker.py}). Each frame, current detections are matched to existing
tracks by solving a Hungarian assignment on the pairwise centroid-distance
matrix; matches beyond \SI{2.0}{\meter} are rejected, unmatched detections start
new tracks, and a track that goes unmatched for more than \num{5} consecutive
frames is retired. This is an offline, centroid-only tracker --- no motion model,
no appearance --- which is adequate here because detection, not association, is
the bottleneck; an IoU-based (SORT-style) tracker is listed as future work.

The tracker also enables a \emph{track-level filter} that improves precision
(Figure~\ref{fig:tracker}).
Each detection carries the classifier's per-frame label; a track is accepted as a
car only if, across its lifetime, it collected at least \num{2} non-rejected
votes, those known votes are a majority of its votes
($\geq \num{0.5}$), and they agree on a single class. A cluster that the
classifier flickers on --- a car in one frame, not-car in the next --- is thus
resolved by temporal majority rather than trusted per frame. Because this filter
consults the whole track, it is an offline post-processing step, consistent with
the offline scope of the system; its contribution is isolated as an ablation row
in Chapter~\ref{ch:evaluation}. The tracker itself is effectively free at
\SI{0.3}{\milli\second\per\frame}.

\begin{figure}[H]
  \centering
  \begin{tikzpicture}[
      font=\footnotesize, node distance=6mm and 14mm,
      det/.style={draw, rounded corners, align=center, fill=blue!5, thick,
        minimum height=10mm, minimum width=20mm, inner sep=3pt},
      dec/.style={draw, rounded corners, align=center, fill=orange!10, thick, inner sep=5pt},
      res/.style={draw, ellipse, align=center, inner sep=3pt, minimum width=22mm},
      arr/.style={-{Stealth[length=2mm]}, thick},
      lbl/.style={font=\scriptsize\itshape, text=black!60, align=center},
    ]
    % one track's per-frame detections, linked by centroid association
    \node[det] (d1) {det @ $t$\\[1pt]{\scriptsize\bfseries car}};
    \node[det, right=of d1] (d2) {det @ $t{+}1$\\[1pt]{\scriptsize\bfseries car}};
    \node[det, right=of d2] (d3) {det @ $t{+}2$\\[1pt]{\scriptsize not-car}};
    \draw[arr] (d1)--(d2);
    \draw[arr] (d2)--(d3);
    \node[lbl, above=1mm of d2] {Hungarian assignment, centroid $\le$\,\SI{2}{\meter}};
    \node[lbl, right=6mm of d3, text width=24mm] (ret)
      {unmatched $>\num{5}$ frames\\$\Rightarrow$ track retired};

    % track-level filter over the whole vote sequence
    \node[dec, below=12mm of d2, text width=70mm] (filt)
      {\textbf{Track-level filter} (whole track, offline)\\[2pt]
       $\ge \num{2}$ known votes \;\, $\wedge$ \;\, known-vote majority $\ge \num{0.5}$
       \;\, $\wedge$ \;\, single agreed class};
    \draw[arr] (d1) -- (filt.north west);
    \draw[arr] (d2) -- (filt.north);
    \draw[arr] (d3) -- (filt.north east);

    \node[res, fill=green!15, below left=9mm and 3mm of filt] (acc) {accept:\\car track};
    \node[res, fill=black!5, below right=9mm and 3mm of filt] (rej) {reject};
    \draw[arr] (filt.south) -- (acc);
    \draw[arr] (filt.south) -- (rej);
  \end{tikzpicture}
  \caption[Tracker and track-level filter]{The tracker and track-level filter. Per-frame detections
    are linked into tracks by Hungarian assignment on centroid distance (matches beyond \SI{2}{\meter}
    rejected; a track unmatched for more than \num{5} frames retired). The track-level filter then
    resolves each track over its \emph{whole} vote sequence rather than per frame: a track is accepted as
    a car only with at least \num{2} known votes, a known-vote majority of at least \num{0.5}, and
    agreement on a single class. In the worked example the votes \{car, car, not-car\} clear all three
    gates, so a classifier flicker in one frame is overruled by temporal majority. Both steps are
    offline; the tracker costs \SI{0.3}{\milli\second\per\frame}.}
  \label{fig:tracker}
\end{figure}

\subsection{The Frozen Pipeline Configuration}
\label{sec:pipeline-config}

Table~\ref{tab:pipeline-config} lists the frozen parameter values that produce
every number in this thesis. They are the committed \texttt{PIPELINE\_CONFIG} and
were not altered to obtain any reported result.

\begin{table}[H]
  \centering
  \caption[Frozen pipeline configuration]{Frozen pipeline configuration (\texttt{src/pipeline.py}). These values
    were fixed before the write-up; no reported number was produced by changing
    them.}
  \label{tab:pipeline-config}
  \begin{tabular}{lll}
    \toprule
    Stage & Parameter & Value \\
    \midrule
    Preprocess & $z$-filter threshold & \SI{-2.0}{\meter} \\
               & voxel size & \SI{0.05}{\meter} \\
               & denoise neighbours / std-ratio & \num{20} / \num{2.0} \\
               & voxel before denoise & true \\
    \midrule
    Ground     & RANSAC iterations & \num{300} \\
               & inlier distance & \SI{0.2}{\meter} \\
               & min.\ normal $z$ & \num{0.5} \\
    \midrule
    Clustering & method & HDBSCAN \\
               & min.\ cluster size / min.\ samples & \num{10} / \num{5} \\
               & clustering voxel size & \SI{0.10}{\meter} \\
    \midrule
    Geometric  & min.\ points & \num{15} \\
               & volume range & $[\num{0.5}, \num{50}]\,\si{\meter\cubed}$ \\
               & max.\ longest extent & \SI{6.0}{\meter} \\
               & max.\ height above ground & \SI{1.5}{\meter} \\
               & max.\ vertical span & \SI{1.8}{\meter} \\
               & max.\ aspect ratio & \num{6.0} \\
    \midrule
    Classifier & checkpoint & \texttt{stage\_b\_scratch\_best.pth} \\
               & input points & \num{512} \\
    \midrule
    Tracker    & max.\ match distance & \SI{2.0}{\meter} \\
               & max.\ disappeared frames & \num{5} \\
               & min.\ track length & \num{2} \\
               & min.\ known votes / ratio & \num{2} / \num{0.5} \\
    \bottomrule
  \end{tabular}
\end{table}

The completion stage that operates on the tracked car detections
(Figure~\ref{fig:pipeline}, right) is the subject of Section~\ref{ch:completion-method}; the evaluation of
the detections themselves is Chapter~\ref{ch:evaluation}.


%% ===== merged from sec_6_completion_method.tex ===== %%

% Point-Cloud Completion: Method (clean-method-first; false starts gathered in a
% closing "Negative results and lessons" subsection). Folds under Chapter 3.
%
% Scope: this section is the completion METHOD and the debugging narrative that
% produced it. The completion RESULTS (donor metric, box metric, held-out
% replication, movers) live in Chapter 4 (Evaluation); here we only state each fix's headline
% effect and defer the full evaluation.
%
% Honesty-phrasing rules applied (personal_agenda P1):
%   - domain-gap failures (#15-19) and Step 1c (#45) presented as FINDINGS, not
%     apologies; #26 inference bug presented as a debugging lesson.
%   - final path = per-car length estimate ON, Step 1c flags OFF.
%   - 23% fallback (#41) stated as a live limitation.
%   - measured vs inferred tiers marked; numbers traceable to findings/artifacts.
%
% Figures (all embedded via \graphicspath from the gitignored output/ tree;
% regenerable from the named scripts):
%   Fig 6.1 fig61_partial_compare.py    -> fig61_partials/       (naive vs KITTI-like partial)
%   Fig 6.2 fig62_box_overlay.py        -> fig62_deblob/box_overlay_deblob_08.png
%           (BEV old vs corrected completion box vs amodal GT; reconstructs the
%            retired PCA+partial-radius normalization, touches nothing frozen)
%   Fig 6.3 viz_completion.py           -> completion_shape_08.png (final output, output/08)
% Floats set [htbp]. Draft-grade renders; may get a cosmetic restyle in the figure pass.
%
\section{Object Completion}
\label{ch:completion-method}

This part of the method takes each detected car and \emph{completes} it:
from the one-sided, partial cluster a single LiDAR frame observes, it predicts
the full car surface. The section presents the finished method in the order a
cluster passes through it --- the network, its training data, the inference-time
normalisation, the input gate, and the length geometry --- and then, separately,
gathers the negative results and lessons that shaped it
(Section~\ref{sec:completion-lessons}). Those false starts are of independent value
--- one is the motivating case for a project-wide rule --- so they are reported as
findings rather than smoothed over, but they are kept out of the method description
so that reads as what the pipeline does, not as how it was arrived at. The scientific
content here is not the network, a standard Point Completion Network (PCN) used off the
shelf, but the data it is trained on and the geometry by which a real partial is
presented to it. The evaluation of the finished method --- whether it adds genuinely
unseen surface, and whether that generalises --- is the subject of
Section~\ref{sec:donor-validation}; here we state each step's headline effect and defer
the rest.

\subsection{Completion Targets: Single-Frame Car Clusters}
\label{sec:completion-what}

Completion runs on the tracked car detections of Section~\ref{sec:meth-detection}. The input is a
\emph{single-frame} cluster in the sensor frame, not the accumulated multi-frame
track. This is a deliberate choice with an evidence basis: accumulating a moving car over many frames produces a
motion smear up to several metres long, which is neither a valid completion
target nor matchable to a car shape. A single frame is one clean, if partial,
observation. The value proposition --- that completing this partial improves the
car's amodal box and recovers surface the frame never saw --- is what Section~\ref{sec:donor-validation}
measures; this section is how the completion is produced.

\subsection{The Completion Network (PCN)}
\label{sec:completion-pcn}

We use PCN (Yuan et al., 2018), an encoder--decoder for point clouds
(\texttt{src/pcn.py}; Figure~\ref{fig:pcn}). The encoder is a stacked PointNet: shared per-point MLPs
lift each point to a local feature, a max-pool summarises the whole cloud into a
global vector, that global vector is concatenated back onto every point, and a
second MLP-plus-max-pool stage produces a single \num{1024}-dimensional shape
embedding. The decoder is coarse-to-fine: a fully-connected stage predicts
\num{1024} coarse seed points, then a \emph{grid-folding} stage attaches a small
2D grid patch to each seed and learns to fold it into local surface geometry,
densifying the coarse skeleton into the final cloud. Training minimises a
two-term Chamfer loss, $\mathrm{CD}_{\text{coarse}} + 0.5\,\mathrm{CD}_{\text{fine}}$;
the fine term is down-weighted so the coarse skeleton is learned first. At
inference the partial is subsampled to \num{256} points before it enters the
encoder.

\begin{figure}[H]
  \centering
  \begin{tikzpicture}[
      font=\footnotesize, node distance=6mm and 7mm,
      blk/.style={draw, rounded corners, align=center, fill=orange!12, thick,
        minimum height=8mm, inner sep=3pt},
      io/.style={draw, ellipse, align=center, fill=black!5, inner sep=2pt},
      arr/.style={-{Stealth[length=2mm]}, thick},
      lbl/.style={font=\footnotesize\itshape, text=black!60},
    ]
    \node[io] (in) {partial\\\num{256} pts};
    \node[blk, right=of in] (ppm) {per-point\\MLP};
    \node[blk, right=of ppm] (mp1) {max-pool,\\concat back};
    \node[blk, right=of mp1] (mp2) {MLP,\\max-pool};
    \node[blk, right=of mp2] (emb) {shape code\\\num{1024}-d};
    \node[blk, below=10mm of emb] (fc) {FC};
    \node[blk, left=of fc] (coarse) {\num{1024}\\coarse seeds};
    \node[blk, left=of coarse] (fold) {grid-folding\\{\scriptsize 2D patch / seed}};
    \node[io, left=of fold] (out) {dense car\\\num{4096} pts};
    \draw[arr] (in)--(ppm); \draw[arr] (ppm)--(mp1); \draw[arr] (mp1)--(mp2);
    \draw[arr] (mp2)--(emb); \draw[arr] (emb)--(fc); \draw[arr] (fc)--(coarse);
    \draw[arr] (coarse)--(fold); \draw[arr] (fold)--(out);
    \node[lbl, above=1mm of mp1] {encoder (stacked PointNet)};
    \node[lbl, below=1mm of coarse] {decoder (coarse-to-fine)};
  \end{tikzpicture}
  \caption[Completion network (PCN)]{The completion network (PCN, \texttt{src/pcn.py}). A stacked-PointNet
    encoder maps the \num{256}-point partial to a \num{1024}-d global shape code;
    a coarse-to-fine decoder predicts \num{1024} seed points and folds a small 2D
    grid patch onto each to produce the dense \num{4096}-point car. Every block is
    learned; the architecture is used off the shelf and unmodified.}
  \label{fig:pcn}
\end{figure}

The architecture is not a contribution and was not modified; the controlled comparison
against PoinTr that justifies keeping PCN --- a small synthetic edge and equivalent
real-data behaviour --- is reported with the completion ablations
(Section~\ref{sec:abl-completion}). What makes PCN
work on real cars is not the weights but the two subsections that follow: the
distribution of partials it is trained on, and the frame a real partial is normalised
into at inference.

\subsection{Training Data: KITTI-Like Partials}
\label{sec:completion-kittilike}

PCN is trained on synthetic partials, and those partials must match the distribution
of the real clusters it meets at inference or the network sees a different input at
test time than it was trained on. A derisk study comparing candidate partial
generators against real seq-08 clusters
(\texttt{scratchpad/derisk\_synth\_vs\_real.py}) isolated three properties a naive
generator gets wrong: it sampled a dense pinhole depth render, skipped the
\SI{0.05}{\meter} voxelisation and ground removal the real input has been through, and
added range-proportional noise (up to \SI{0.15}{\meter} at \SI{30}{\meter}) that
washed out the scan-ring banding real Velodyne data shows.

The generator used (\texttt{\_render\_kitti\_like} in \texttt{src/train\_pcn.py}) makes the synthetic partial match the real one: it orients the
ShapeNet car to gravity, casts a single HDL-64E viewpoint (64 elevation beams,
\SI{0.09}{\degree} horizontal resolution) from a sensor at \SI{1.73}{\meter} and
a sampled range of \SIrange{8}{30}{\meter}, adds small \emph{absolute} noise
(\SI{0.015}{\meter}), voxelises at \SI{0.05}{\meter} to match the pipeline, and
removes points within \SI{0.30}{\meter} of the ground plane. The dense complete
mesh sample remains the target. PCN retrained on these partials
(\texttt{pcn\_kitti\_best.pth}, best val loss \num{0.1246}) is in-distribution
with the real input, and on a held-out synthetic set produces clean cars, not
blobs (Chamfer \SI{0.16}{\meter}, F-score@\SI{0.1}{\meter} \num{0.76}).
Figure~\ref{fig:kitti-like-partial} contrasts the naive and KITTI-like partials
for the same ShapeNet car.

\paragraph{Training.} PCN is trained on the KITTI-like synthetic partials (input
\num{256} points, target the dense mesh sample) for \num{100} epochs with Adam
(learning rate \num{e-4}, weight decay \num{e-6}) and a step schedule that halves
the rate every \num{40} epochs, at batch size \num{8}, minimising the two-term
Chamfer loss above. No real cars enter training; the reported checkpoint
(\texttt{pcn\_kitti\_best.pth}) is the best-validation-loss epoch (\num{0.1246}).
The subsequent geometry work is all inference-time and does not retrain the
network.

\begin{figure}[htbp]
  \centering
  \includegraphics[width=\linewidth]{fig61_partials/partials_2668f2d9b8a2daad95551e0d9e50cb0d.png}
  \caption[Training-input distribution]{The training-input distribution: one ShapeNet car in
    three views, with its complete-surface target (left, grey) and two candidate
    training \emph{inputs}. The naive pinhole partial (centre, orange) is dense and
    uniform; the KITTI-like partial (right, blue) reproduces what real LiDAR
    delivers --- scan-ring banding, single-viewpoint coverage, \SI{0.05}{\meter}
    voxelisation, and the ground cut. Neither is a completion output; ``better''
    here means closer to the real input distribution, not more complete. Training on
    the KITTI-like distribution is what keeps the network in-distribution on real data.}
  \label{fig:kitti-like-partial}
\end{figure}

\subsection{Inference-Time Partial Normalisation}
\label{sec:completion-normalisation}

PCN is trained on shapes expressed in a fixed canonical car frame, centred on the
true car centre and scaled by the true car radius, so at inference a real partial
must be mapped into that \emph{same} frame before it is completed --- anything else
hands the network a normalisation it never saw during training. Reconstructing that
frame from a one-sided partial is non-trivial, because a partial reveals neither the
true centre nor the true radius directly. The inference path
(\texttt{estimate\_canonical\_frame}) reconstructs it in three steps:

\begin{itemize}
  \item \textbf{Reorientation.} The sensor frame is rotated into the canonical
    car frame (X${}={}$width, Y${}={}$up, Z${}={}$length): gravity maps to the
    up axis, and the horizontal heading axis maps to length. Heading comes from the
    L-shape footprint fit (Section~\ref{sec:completion-lshape}), not a principal-axis
    alignment.
  \item \textbf{Full-car centre estimate.} A one-sided partial's bounding-box
    centre is not the car centre, so the centre is corrected on each axis: shift
    \emph{down} to undo the upward bias from the ground cut (the dominant
    correction, \SI{0.25}{\meter}); push \emph{outward} along width toward the
    occluded far side to a half-car-width prior (\SI{1.9}{\meter}); and push
    along length toward the occluded far end (Section~\ref{sec:completion-length}).
  \item \textbf{Scale.} Normalise by the maximum point-to-centre distance divided
    by a fixed scale correction (\num{1.137}), the calibrated ratio between a
    one-sided partial radius and the true car radius.
\end{itemize}

With this path the same real seq-08 clusters recover the correct pose and scale; the
ablation isolating this normalisation against a naive partial-frame alternative, and the
box-overlap gains it produces, is reported with the completion ablations
(Section~\ref{sec:abl-completion}). Viewed in the
canonical car frame, the corrected completions are recognisably car-shaped: they
recover the car footprint in bird's-eye view and a rough side silhouette
(roofline included) even from very sparse input --- a clean car footprint from a
\num{64}-point cluster in the extreme case. We are nonetheless careful about the
strength of the claim: crispness varies with input density (cleaner on
well-observed cars, blobbier on the sparsest), and a completion is a filled
\num{4096}-point cloud rather than a CAD-clean surface. The sharpest completions
are the \emph{in-distribution synthetic} ones of
Section~\ref{sec:completion-kittilike} (Chamfer \SI{0.16}{\meter}); on real data
the normalisation's value is corrected geometry --- pose, scale, and the
surface-coverage gains quantified in Section~\ref{sec:donor-validation}. A calibration study on
synthetic cars further identified the residual dominant error as
\emph{centre estimation}, not scale --- which is what motivates the geometry work
in the rest of this section.

\subsection{L-Shape Gating of Completion Inputs}
\label{sec:completion-lshape}

Even with correct normalisation, some clusters cannot complete into a plausible
car because the \emph{input} is not one car: HDBSCAN fragments (part of a car) and
merges (two cars, or a \SI{90}{\degree} heading error) both feed the completion
network implausible input. To detect them we fit the cluster's ground-plane footprint with a
search-based L-shape fit (Zhang et al., 2017), which scores edge adherence over
candidate yaw angles and stays well-defined on near-square footprints where PCA
degenerates. The fitted footprint gives an input-quality gate: a
cluster is rejected as a fragment if its fitted length is below \SI{2.7}{\meter},
or as a merge if its fitted width exceeds \SI{2.3}{\meter}.

On seq-08 (\texttt{scratchpad/ab\_heading.py}) \emph{every} implausible completion
originated from a fragment or merge input, so the gate's effect on completion quality
is a clean ablation, reported in Section~\ref{sec:abl-completion}. An honest secondary
finding from the same study: the L-shape heading itself is neutral against plain PCA on
real data --- the fit earns its place by \emph{detecting bad inputs}, not by giving a
better heading.

\subsection{Per-Car Length Estimation}
\label{sec:completion-length}

The canonical-frame centre estimate re-centres the occluded lateral side and the
ground cut, but a partial truncated at the occluded far \emph{end} is normalised
around its observed near portion, so PCN's output stops short of the far end. The
donor metric of Section~\ref{sec:donor-validation} confirmed the far end as the single worst-covered
region. Two successive fixes address it, both inference-only (no retraining).

\paragraph{The fixed length prior.} A longitudinal, extend-only
push mirrors the width prior along the length axis: the centre is pushed toward
the occluded far end so the completion reaches an assumed full car length of
\SI{4.14}{\meter} (the real amodal-GT median). This more than doubled far-end
coverage on the donor metric and reversed the box-length regression (signed
$\Delta L$ \SI{-0.44}{}$\to$\SI{-0.32}{\meter}). A larger \SI{4.5}{\meter} prior
was rejected because it broke the hallucination guard by over-extending genuinely
compact cars.

\paragraph{The per-car estimate.} A single \SI{4.14}{\meter}
prior is a population median, so it is wrong in both tails --- it over-extends
compact cars and under-extends long ones. It is replaced by a per-car length,
\texttt{track\_length\_estimate}: the \nth{90} percentile of the car's own
gate-passed footprint-length fits, plus a \SI{0.12}{\meter} bias correction,
falling back to the fixed prior when a track has fewer than \num{5} gate-passed
frames. The estimator was chosen offline against amodal GT with no PCN inference:
a length-from-width shortcut is impossible (the correlation between true length
and width is \num{0.018}), and the aggregate must be a \emph{quantile}, not the
maximum, because a single contaminated L-shape fit sets the maximum. This is a
per-track quantity, so the pipeline (which owns the track) computes it and passes
it in; the completion module itself stays stateless. It removes the compact
over-extension (compact signed $\Delta L$ \SI{+0.295}{}$\to$\SI{+0.063}{\meter})
and improves the box metric overall (BEV IoU \num{0.747}$\to$\num{0.771}), at a
measured cost in pooled donor coverage that Section~\ref{sec:donor-validation} reports in full.

\paragraph{Limitation.} The fallback is not rare: on the
regenerated seq-08 output, \num{119} of \num{518} completed tracks
(\SI{23}{\percent}) fell back to the fixed prior because they had fewer than five
gate-passed frames. This is above a comfortable threshold and is stated as a live
limitation, not a solved problem; it is revisited in the discussion.

\subsection{The End-to-End Completion Path}
\label{sec:completion-summary}

To summarise, a single-frame car cluster is completed by: fitting its footprint
and rejecting fragments and merges (L-shape gate); reorienting it into the
canonical car frame with no PCA; estimating the full-car centre (down-shift for
the ground cut, width prior for the occluded side, per-car length estimate for
the occluded far end); normalising by the scale-corrected radius; running the
KITTI-like-trained PCN; and rotating the output back to the sensor frame. Every
step is inference-time geometry except the one-off retraining on KITTI-like
partials; no step fine-tunes on the evaluation cars.

Figure~\ref{fig:completion-shape} shows what this path produces on real seq-08
cars across the input-density range, in the canonical car frame. The completed
cloud (green) fills the car footprint and a rough side silhouette that the
single-frame partial (blue) leaves one-sided --- recognisably car-shaped down to a
\num{64}-point input, with fitted dimensions in the expected car range
($L \approx \SI{4.2}{\meter}$, $W \approx \SI{1.9}{\meter}$,
$H \approx \SI{1.6}{\meter}$ on the well-observed cars). Consistent with the
caveats above, the sparsest inputs complete more coarsely and can inflate the
width (Section~\ref{sec:donor-validation}), and the clouds are filled rather than CAD-clean;
the point here is qualitative shape, with the quantitative real-data assessment ---
whether the completion adds genuinely unseen surface, and whether that survives a
pre-registered held-out replication --- being what Section~\ref{sec:donor-validation} sets out to measure.

\begin{figure}[htbp]
  \centering
  \includegraphics[width=\linewidth]{completion_shape_08.png}
  \caption[Completions on real seq-08 cars]{What the final completion path produces on real seq-08 cars, four
    tracks spanning the input-density range (\num{4679} down to \num{64} input
    points), each in the canonical car frame: bird's-eye view (left of each pair)
    and side elevation (right). Blue is the single-frame partial input, green the
    PCN completion. The completion recovers the full car footprint and a rough
    roofline silhouette from a one-sided partial; crispness and dimensional
    accuracy degrade on the sparsest inputs (the \num{64}-point car shows the
    width inflation discussed in Section~\ref{sec:donor-validation}). Rendered by
    \texttt{scratchpad/viz\_completion.py} from the final-config \texttt{output/08}.}
  \label{fig:completion-shape}
\end{figure}

\subsection{Negative Results and Lessons}
\label{sec:completion-lessons}

The method above was not reached in one step. Two false starts shaped it, and each
carries a lesson worth stating in its own right rather than smoothing over. They are
gathered here so that the method description reads as what the pipeline does. A third
intervention that was tried and rejected --- a length-dependent far-end
compensation --- is reported as a completion ablation
(Section~\ref{sec:abl-completion}).

\paragraph{The blobs were not a synthetic-data problem.}
\label{sec:completion-domaingap}
The first attempts to run PCN on real cars all produced ``blobs'' --- rounded,
car-sized clouds with no recognisable car structure. This was initially attributed to
the training data being synthetic, and the response was a series of sim-to-real
remedies: LiDAR-noise augmentation of the synthetic partials, and real-data
fine-tuning. None of them fixed the blobs, and the fine-tuning attempt is recorded as
a negative result --- on this data, fine-tuning PCN on mined real pairs did not recover
real-car geometry. The blobs had two causes, and neither was the network weights: the
training-input distribution (Section~\ref{sec:completion-kittilike}) and the inference
normalisation (Section~\ref{sec:completion-normalisation}), the two method steps above.
The lesson is to suspect the data and the input pipeline before the model.

\paragraph{Train and inference semantics must match.}
\label{sec:completion-inferencebug}
Even after retraining on KITTI-like partials, real clusters still blobbed, and the
cause is the sharper lesson. An earlier inference path normalised each real partial by
a 3D PCA frame and the \emph{partial's} own centroid and radius --- a normalisation the
model had never seen, since it was trained in the canonical car frame of
Section~\ref{sec:completion-normalisation}. This silent train/inference mismatch alone
produced blobs, and it did so for \emph{every} checkpoint, including in-distribution
synthetic ones (it worsened Chamfer by roughly \num{3.5}$\times$ even where the model
is otherwise clean). Because the two frames each looked internally reasonable, the
divergence was invisible until measured. This incident is the motivating example for
the project's rule that any external integration or reuse gets a written semantics map
--- coordinate frame, centring, and scale conventions --- reconciled before code is
run.

\section{The Donor-Frame Completion Metric}
\label{sec:donor-metric}

Evaluating the completion output requires a metric, and on real LiDAR there is no
complete ground-truth surface to score against. This section explains why the
obvious choice --- Chamfer distance to a pseudo ground truth --- is invalid, then
defines the \emph{donor-frame} metric used in its place. The metric is
validated and applied to the completion output in Section~\ref{sec:donor-validation};
it is defined on static cars by construction, with moving cars treated separately in
Section~\ref{sec:results-movers}.

\subsection{The Invalidity of Chamfer Distance on Real Cars}
\label{sec:chamfer-invalid}

The natural way to score a completed car is Chamfer distance (CD) against a
ground-truth surface. Real LiDAR offers no such surface: a car is seen once, from
one side. The obvious substitute is a \emph{pseudo}-ground-truth cloud --- the
labelled points of one static car accumulated across every ego frame that
observed it. This subsection explains why that substitute is not a valid target,
and hence why the completion work of Section~\ref{ch:completion-method} needed a different metric
altogether.

The fundamental limitation is that the accumulation is itself one-sided. A parked car is
observed from a road running past one of its flanks, so every frame sees broadly
the same near faces --- the road-facing side, the near end, the roofline edge ---
and none sees the occluded far side or the underbody. Pooling those frames
densifies the \emph{visible} surface but adds almost no genuinely new surface: the
pseudo-ground-truth is a denser partial, not a whole car.

Scored against such a reference, plain Chamfer \emph{rewards under-completion}.
The distance is symmetric, so every predicted point pays the distance to its
nearest reference point. A completion that correctly reconstructs the occluded
far side places points where the reference has nothing, and each such point adds
cost; a model that predicts only the already-observed surface pays nothing for
what it omits. The raw partial --- which by construction sits exactly on observed
surface and predicts nothing new --- therefore scores the \emph{lowest} (best) CD
of any cloud, on every real example. A metric that ranks the
un-completed input first cannot be used to measure completion.

This is a measured failure rather than a conjecture: across the real seq-08 cars
we scored, the raw partial had the lowest pseudo-ground-truth CD of the raw /
mirrored / completed clouds in every case (renders in
\texttt{output/experiments/verify\_pcn\_step2/}). Its consequence is that when the geometry work
began, completion had no valid real-data metric at all --- the in-distribution
synthetic CD of Section~\ref{ch:completion-method} (\SI{0.16}{\meter}) measures the network on synthetic
cars but says nothing about a real cluster. The next subsection turns the very
one-sidedness that breaks Chamfer into the basis of a metric that works.

\subsection{The Donor-Frame Occluded-Side Principle}
\label{sec:donor-principle}

The idea turns the one-sidedness of LiDAR from a problem into the measurement.
For a static car observed across many frames, we complete it from \emph{one}
frame's pipeline cluster and then score the completed cloud only against surface
that the \emph{other} frames --- the \emph{donor} frames --- saw but the input
frame did not (Figure~\ref{fig:donor-principle}). Two properties make this a fair
test of completion:

\begin{itemize}
  \item The raw partial covers \emph{none} of that donor-only surface by
    construction, so any coverage the completion achieves is genuinely added,
    previously-unseen surface --- not a rescoring of the input points.
  \item The reference is built from ground-truth labelled points only, never
    from a pipeline output, so no detection or completion error can leak into
    the thing we measure against. The metric is a \emph{paired} comparison of
    raw versus completed on identical inputs, so it tolerates reference noise
    the same way the box evaluation does.
\end{itemize}

\begin{figure}[htbp]
  \centering
  \includegraphics[width=0.86\linewidth]{donor_principle.png}
  \caption[Donor-frame occluded-side principle]{The donor-frame occluded-side principle (schematic; a car seen from
    above, the road below). The input frame observes the near flank (blue); as the
    ego drives past, the \emph{donor} frames additionally observe the two ends
    (amber), which together form the \emph{novel set} --- the surface the metric
    scores against. The far flank (grey) is occluded from every frame and is not
    part of the reference. Because the novel set lies at least $\tau$ from every
    input point, the raw partial covers none of it by construction. Illustrative
    geometry, not measured data; real per-car coverage on seq~08 is
    Figure~\ref{fig:donor-coverage}.}
  \label{fig:donor-principle}
\end{figure}

\paragraph{Relation to existing completion metrics.} Completion on real LiDAR has
been scored before without a complete target, but not against the previously-unseen
surface specifically. The KITTI completion metrics introduced with
PCN~\cite{yuan2019pcnpointcompletionnetwork} --- Fidelity, Minimal Matching Distance, and
Consistency --- reference the input, a synthetic car prior, and the model's own cross-frame
outputs respectively, so a completion that merely reproduces its input can score well on
all three. Ren et al.~\cite{ren2022selfsupervised} come closest to the construction used
here: they too accumulate each tracked vehicle across frames to form a pseudo ground truth.
The difference is what the accumulation is scored against. They compare the completed cloud
to the accumulation \emph{in full}, mixing the near surface the input already observed with
the small amount the other views add; the donor metric keeps only the novel set --- the
surface at least $\tau$ from every input point --- so already-observed surface contributes
nothing and the raw partial scores zero by construction. This is also what separates the
metric from semantic scene completion~\cite{Behley_2019_ICCV}, which scores voxel occupancy
over a whole scene rather than surface coverage of one object. The novel-set restriction,
the mirrored baseline of Section~\ref{sec:donor-definition}, and the per-band hallucination
guard of Section~\ref{sec:donor-guard-lesson} are what make the metric measure
occluded-surface recovery specifically, rather than overall resemblance to an accumulated
cloud.

\subsection{Definition of the Donor-Frame Coverage Metric}
\label{sec:donor-definition}

For each (frame, instance) true-positive pair on the well-observed static cars,
we form three point sets and score them against a common target.

\paragraph{Inputs.} The \textbf{input} is the production pipeline's cluster for
that car in that frame, completed through the production path (the KITTI-like
PCN checkpoint behind the L-shape input gate). The \textbf{donor cloud} is the
union of the instance's ground-truth (labelled-car) points from every observing
frame \emph{except} the input frame, transformed to the world frame and
voxelised at \SI{3}{\centi\meter}.

\paragraph{Novel set (the visibility mask).} From the donor cloud we keep only
points that lie at least $\tau$ from \emph{every} input point (nearest-neighbour
test). This \emph{novel set} is exactly the surface the input frame did not
observe. We use $\tau = \SI{0.15}{\meter}$ as the primary radius and
\SI{0.10}{} / \SI{0.20}{\meter} as stability checks.

\paragraph{Scored clouds.} We score three clouds against the novel set so that
completion is judged against meaningful baselines rather than in isolation:
\begin{itemize}
  \item \textbf{raw} --- the partial input (the floor: covers no novel surface
    by construction);
  \item \textbf{mirrored} --- the input unioned with its own reflection across
    the car's estimated symmetry plane (a cheap, hallucination-free way to add
    surface; the baseline completion must beat to justify a learned model);
  \item \textbf{completed} --- the PCN output.
\end{itemize}

\paragraph{Scores.} Each cloud gets three numbers, all one-directional from the
novel set to the cloud:
\begin{itemize}
  \item \textbf{coverage} (\texttt{cov@0.1}) --- the fraction of novel points
    within \SI{0.1}{\meter} of the cloud (the primary quality score: how much
    unseen surface was recovered);
  \item \textbf{median novel distance} --- the median distance from novel points
    to the cloud (a threshold-free companion to coverage);
  \item \textbf{out-of-box} --- the fraction of the cloud's points falling
    outside the amodal ground-truth box (plus a \SI{0.2}{\meter} margin): a
    \emph{hallucination guard}, since a model can trivially raise coverage by
    inventing surface everywhere.
\end{itemize}

\paragraph{Eligibility.} The metric applies to any completion whose input is a
\emph{single, well-formed object cluster} --- one cluster corresponding to one
car, with a footprint consistent with a whole vehicle --- so that there is a
coherent object shape to score. In this pipeline that admission is realised by the
L-shape input gate of Section~\ref{ch:completion-method} (fitted footprint length
$\geq \SI{2.7}{\meter}$ and width $\leq \SI{2.3}{\meter}$, rejecting fragments and
merges), but the metric itself does not depend on that particular gate: any
admission rule that isolates a single object serves in its place.

\paragraph{Statistics.} A pair counts if its input is admitted and its novel set
has at least \num{100} points at the primary $\tau$. Because the frames of one
parked car are autocorrelated, we reduce each car to its per-car median and test
across cars with the Wilcoxon signed-rank test --- the same paired, per-car design
used for the box evaluation.

\subsection{Summary and Methodological Limitations}
\label{sec:meth-limitations}

This chapter has set out the full method: the modular detection pipeline that turns a
raw scan into car detections, the single-frame completion path with its input gate and
length prior, and the definition of the donor-frame coverage metric. Several limitations
are built into these choices and are flagged here so the evaluation is read with them in
view; each is taken up in full in the discussion of Chapter~\ref{ch:conclusion}.

\begin{itemize}
  \item \textbf{Ground model.} A single global RANSAC plane assumes near-planar
    ground; it holds on the flat urban sequences used here but would degrade on
    strong slopes, where a grid-based model (future work) would be needed
    (Section~\ref{sec:pipeline-ground}).
  \item \textbf{Offline by design.} The tracker and track-level filter consult the
    whole sequence, and at $\approx$\SI{1.6}{\fps} the pipeline is far below sensor
    rate; real-time operation is out of scope (Section~\ref{sec:pipeline-overview}).
  \item \textbf{Single-frame, static-car completion.} Completion runs on one frame's
    cluster, and the donor metric is defined on static cars; moving cars are treated
    separately and are not scored for completion accuracy
    (Section~\ref{sec:completion-what}).
  \item \textbf{Length-prior fallback.} The per-car length estimate falls back to the
    fixed \SI{4.14}{\meter} prior on \SI{23}{\percent} of tracks (those with fewer
    than five gate-passed frames), and a residual far-end under-extension on long cars
    is a characterised, not eliminated, limit (Sections~\ref{sec:completion-length}
    and~\ref{sec:completion-step1c}).
  \item \textbf{Synthetic-only completion training.} PCN sees only synthetic
    KITTI-like partials; real-data fine-tuning was a negative result, so a sim-to-real
    residual remains (Section~\ref{sec:completion-domaingap}).
  \item \textbf{Metric admissibility.} The donor metric requires a single well-formed
    object cluster and multi-frame donor observations, so it does not score fragments,
    merges, or movers (Section~\ref{sec:donor-definition}).
\end{itemize}

With these caveats stated, Chapter~\ref{ch:evaluation} puts each component to the test
--- the detection operating point and its ablations, the recall root-cause, and the
donor metric's validation and completion results --- and Chapter~\ref{ch:conclusion}
returns to the limitations above alongside the evaluation-validity questions the results
raise.

